% ============================================
% Johns Hopkins Letter Template
% Assistant Professor Cover Letter – Colby College
% ============================================

\documentclass[11pt,letterpaper]{letter}

\usepackage{graphicx}
\usepackage{eso-pic}
\usepackage{geometry}
\usepackage{microtype}
\usepackage[hidelinks]{hyperref}

% ----- Page Setup -----
\geometry{left=25mm,right=25mm,top=25mm,bottom=25mm}

% ----- Logo Placement (foreground, first page only) -----
\newcommand{\PutJHULogoFG}[1]{%
  \AddToShipoutPictureFG*{%
    \AtPageUpperLeft{%
      \hspace{10mm}%
      \raisebox{-38mm}{%
        \includegraphics[width=6cm]{#1}%
      }%
    }%
  }%
}

% ----- Sender Information -----
\signature{Decory Edwards}
\address{Department of Economics \\ Johns Hopkins University \\ Baltimore, MD 21218 \\ \texttt{dedwar65@jh.edu}}

\begin{document}
\PutJHULogoFG{Images/jhulogo.png}

\begin{letter}{Search Committee \\
Department of Economics \\
Colby College \\
4000 Mayflower Hill \\ 
Waterville, ME 04901 \\ USA}

\opening{Dear Members of the Search Committee,}

I am writing to express my interest in the tenure-track Assistant Professor position in Macroeconomics in the Department of Economics at Colby College, beginning Fall 2026. I am a Ph.D. candidate in Economics at Johns Hopkins University, advised by Professors Christopher D. Carroll, Nicholas W. Papageorge, and Stelios Fourakis, and I expect to receive my degree in May 2026. My research fields are macroeconomics, wealth and income dynamics, and computational economics.

My research agenda is unified by a focus on understanding the determinants of wealth inequality and the mechanisms through which heterogeneity in returns to wealth shapes aggregate outcomes. In my job-market paper, I estimate the degree of heterogeneity in individual portfolio returns using micro-data from the Survey of Consumer Finances and simulate a structural model in which households face idiosyncratic return risk. I show that allowing for persistent heterogeneity in returns substantially improves the model’s ability to match the empirical distribution of wealth, offering a tractable way to connect micro-level return dispersion to macroeconomic inequality.

In a second project, I use the Health and Retirement Study to study the relationship between interpersonal trust and portfolio performance. Building on the literature that finds a hump-shaped relationship between trust and income, I show that a similar relationship holds for individual returns to wealth measured in the HRS data, highlighting a behavioral channel through which social attitudes shape saving and investment outcomes. This work also provides new estimates of the persistent component of returns to net worth, which motivate the calibration of heterogeneity in my structural model.

Colby College’s identity as a liberal arts institution committed to undergraduate teaching, close faculty-student interaction, and robust research support is an excellent fit for my goals. The Department of Economics emphasizes teaching excellence, inclusion, and working with highly engaged undergraduates in the classroom and research settings. I am particularly drawn to the opportunity to contribute to the macroeconomics sequence at both the introductory and intermediate undergraduate levels, while also offering elective courses shaped by my scholarship—such as \textit{Wealth Dynamics and Macro Risk} or \textit{Computational Macroeconomics \& Policy}. The 4.5-course annual teaching load at Colby affords faculty the ability to remain research-active while engaging deeply in the undergraduate liberal arts environment. 

\textbf{Teaching.} My teaching experience centers on undergraduate macroeconomics and an upper-level macro policy seminar, where I have led weekly discussion sections and held regular office hours for classes ranging from 16 to 40 students. My approach is student-centered: I aim to help students “convince themselves” that they have moved from confusion to understanding by holding structured office-hour coaching that builds trust and confidence. At Colby I am prepared to teach \textit{Principles of Macroeconomics}, \textit{Intermediate Macroeconomic Theory}, \textit{Money \& Banking}, and electives in quantitative macroeconomics. I would also welcome designing and leading a course such as \textit{Computational Macroeconomics and Public Policy} emphasizing reproducible workflows, data analysis, and evidence-based decision-making.

I believe I would be a strong fit because my computational and empirical toolkit allows me to (i) produce robust, data-backed estimates of wealth returns and distributional outcomes, (ii) build and validate models that speak to macroeconomic issues in the broad liberal‐arts teaching context, and (iii) collaborate with faculty and undergraduates to present insights in accessible and actionable formats. I am excited by the possibility of contributing to Colby’s mission of combining scholarly excellence with teaching and mentoring of engaged undergraduates.

I would be honored to join Colby College’s Department of Economics and contribute to its teaching, research, and mentoring mission. Thank you for your consideration; I welcome the opportunity to discuss my fit with the Department at your convenience.

\closing{Sincerely,}

\end{letter}
\end{document}
